\documentclass[10pt]{article}

\usepackage[margin=0.8in]{geometry}
\usepackage{amsmath,amssymb,mathtools}
\usepackage{booktabs}
\usepackage{tabularx}
\usepackage{algorithm}
\usepackage{algpseudocode}
\usepackage{fancyvrb}
\usepackage{hyperref}

\newcommand{\Var}{\mathrm{Var}}
\newcommand{\Cov}{\mathrm{Cov}}
\DeclarePairedDelimiter{\abs}{\lvert}{\rvert}
\DeclarePairedDelimiter{\norm}{\lVert}{\rVert}
\newcommand{\ket}[1]{\lvert #1 \rangle}
\newcommand{\bra}[1]{\langle #1 \rvert}
\newcommand{\expval}[1]{\langle #1 \rangle}

\title{Joint Cluster-Partition and Orbital-Basis Optimization\\
\large A technical note on \texttt{cluster\_number\_decomposition\_optimization.py}}
\date{\today}
\author{}

\begin{document}
\maketitle
\vspace{-2em}

\section{Introduction}

Consider a molecular Hamiltonian $H$ in a basis of $\mathrm{norb}$ spatial orbitals and a
fixed-particle-number state $\ket{\psi}$ (typically a DMRG ground state). A cluster number
decomposition partitions the orbitals into clusters $C_1,\dots,C_k$, each with total occupation
operator $N_{C_i}=\sum_{p\in C_i}n_p$; if $\ket{\psi}$ is close to a joint eigenstate of every
$N_{C_i}$, each acts as an approximate quasisymmetry, enabling the Hilbert-space truncations
this codebase is built around. A decomposition is a pair: the partition $P$, and the orbital
basis $U$ in which the clusters are defined, since rotating the orbitals changes which
operators the $N_{C_i}$ actually are.

The sister module \texttt{cluster\_numbers\_metrics.py} fixes $P$ by hand and only optimizes
$U$. This module, documented here, searches over both jointly. Exhaustive search over $P$ is
impossible -- partitioning $\mathrm{norb}$ orbitals has $B_{\mathrm{norb}}$ (Bell number)
possibilities, already $B_{16}\approx10^{10}$ for a modest active space -- so it uses a
heuristic \emph{beam search} instead: build finer partitions one split at a time, keep a
bounded set of promising candidates, and re-optimize $U$ locally after each split.

\section{Objects and Notation}

\paragraph{Decomposition.} A \emph{decomposition} is a pair $(P,U)$: $P$ a partition of
$\{0,\dots,\mathrm{norb}-1\}$ into clusters, and $U$ the $\mathrm{norb}\times\mathrm{norb}$
orbital rotation relative to a fixed reference basis (molecular orbitals, MOs, or natural
orbitals, NatOs). $U$ is parameterized by a flat vector $x\in\mathbb{R}^{\binom{\mathrm{norb}}{2}}$:
$x$ fills the strict upper triangle of an antisymmetric generator $A$ (in
\texttt{np.triu\_indices(norb,k=1)} order, row-major over $p<q$), and $U=\exp(A)$. Since $A$ is
real and antisymmetric, $U$ is always a \emph{proper rotation}, $U\in SO(\mathrm{norb})$:
$U^\top U=I$ and $\det U=+1$, never merely orthogonal, never a generic complex unitary. This is
the one rotation convention used throughout this module and in the routines it reuses from
\texttt{src/cluster\_number\_operators.py}.

\paragraph{RDM data.} All cost functions are built from reduced density matrices (RDMs) of
$\ket{\psi}$ in the current rotated basis, bundled in an \texttt{RDMData} object:

\begin{center}
\begin{tabular}{ll}
\toprule
Field & Content \\
\midrule
\texttt{D} & spin-summed 1-RDM (always needed) \\
\texttt{Gamma} & spin-summed 2-RDM (always needed) \\
\texttt{rdm3}, \texttt{rdm4} & spin-summed 3-/4-RDM (commutator cost only) \\
\texttt{h1e}, \texttt{g2e\_full} & one-/two-electron integrals (commutator cost only) \\
\bottomrule
\end{tabular}
\end{center}

\paragraph{Cost functional.} Every cost function has the same shape: a sum over clusters of a
per-cluster cost built from that cluster's number-operator statistics,
$C(P,U)=\sum_{c\in P} c(N_c)$, evaluated at rotation $U$. Section~\ref{sec:costs} lists the five
supported cost functions.

\section{Cost Functions}
\label{sec:costs}

\begin{table}[h]
\centering
\begin{tabularx}{\textwidth}{lXll}
\toprule
Name & Formula & Needs & Cost \\
\midrule
variance & $\sum_C \Var(N_C)^p$ & 1-,2-RDM & $O(\mathrm{norb}^5)$ \\
extremality & $\sum_C \expval{N_C}(2\abs{C}-\expval{N_C})$ & 1-RDM & $O(\mathrm{norb}^2)$ \\
eval\_eq & $\sum_C\left[\Var(N_C)+(\expval{N_C}-e_C)^2\right]$ & 1-,2-RDM & $O(\mathrm{norb}^5)$ \\
mixed & $c_1\!\cdot\!\text{variance}+c_2\!\cdot\!\text{extremality}$ & 1-,2-RDM & $O(\mathrm{norb}^5)$ \\
commutator & $\sum_C \expval{[H,N_C]^2}$ & 1-,2-,3-,4-RDM, $h_{1e}$, $g_{2e}$ & see text \\
\bottomrule
\end{tabularx}
\caption{The five cost functions built into this module. $\expval{\cdot}=\bra{\psi}\cdot\ket{\psi}$
throughout. $p$ is a tunable exponent (default 1); $e_C$ are target integer eigenvalues.}
\end{table}

The variance exponent $p$ raises each cluster's variance individually before summing, not the
total. The extremality cost needs only the 1-RDM: it is a downward parabola in $\expval{N_C}$,
zero at $\expval{N_C}\in\{0,2\abs{C}\}$ and maximal at half filling; minimizing it pushes every
cluster's occupation toward an integer extreme.

\paragraph{The commutator cost.} $F(C)=\expval{[H,N_C]^2}$ measures how far $N_C$ is from
commuting with $H$ on $\ket{\psi}$: $F(C)=0$ means $N_C$ is an exact symmetry generator on this
state. Since $H$ and $N_C$ are Hermitian, $[H,N_C]$ is anti-Hermitian, so
$F(C)=-\norm{[H,N_C]\ket{\psi}}^2\le0$ always; the code returns $-F(C)=\norm{[H,N_C]\ket{\psi}}^2\ge0$,
so that, like the other four costs, smaller is better. Expanding $[H,N_C]$ (a 1-/2-body
operator) against another 1-body operator keeps it a 1-/2-body operator, so $F(C)$ contracts
the 1-,2-,3-,4-RDM with $h_{1e}$ and $g_{2e}$; the full formula is in
\texttt{src/cluster\_number\_operators.py}, not reproduced here.
Two implementations exist. \texttt{number\_commutator\_cost\_v1} recomputes an
$O(\mathrm{norb}^8)$-class tensor contraction independently for every cluster.
\texttt{number\_commutator\_cost\_v2} observes that every term of $F(C)$ is linear in the
cluster indicator $I_C\in\{0,1\}^\mathrm{norb}$, used exactly twice, so $F(C)=I_C^\top K I_C$ is
a quadratic form for a single $\mathrm{norb}\times\mathrm{norb}$ kernel $K$ depending only on
$H,\ket{\psi}$, not on $C$; $K$ is built once per rotation (the same $O(\mathrm{norb}^8)$-class
cost, but paid only once), and every cluster's score is then a cheap $O(\mathrm{norb}^2)$
bilinear form. Because it needs the 3-/4-RDM at \emph{every} candidate rotation, it is too
expensive to drive the search itself -- Sections~\ref{sec:fiedler}--\ref{sec:beam} always use
the cheap variance cost instead -- and is reserved for the single polish step
(Section~\ref{sec:polish}) on a chosen trajectory endpoint.

\section{Cluster Splitting: Fiedler Bisection with a Minimum-Size Floor}
\label{sec:fiedler}

\begin{center}
\begin{tabular}{ll}
\toprule
Symbol & Meaning \\
\midrule
$C$, $k=\abs{C}$ & the cluster being split, and its size \\
$n_1,n_2$ & one-/two-body number statistics restricted to $C$ under the current $U$ \\
$W_{pq}$ & $\abs{\Cov(n_p,n_q)}=\abs{n_2[p,q]-n_1[p]n_1[q]}$, $p\ne q \in C$, $W_{pp}=0$ \\
$L$ & graph Laplacian of $W$: $L=\mathrm{diag}(W\mathbf{1})-W$ \\
$v$ & Fiedler vector: eigenvector of $L$ for its second-smallest eigenvalue \\
$k_{\min}$ & \texttt{min\_child\_cluster\_size}: floor on each resulting child's size \\
\bottomrule
\end{tabular}
\end{center}

\begin{algorithm}[h]
\caption{\textsc{FiedlerBipartition}$(n_1,n_2,C,k_{\min})$}
\label{alg:fiedler}
\begin{algorithmic}[1]
\small
\If{$k<2$} \textbf{raise} error: nothing to split \EndIf
\If{$k<2k_{\min}$} \textbf{raise} error: no split of $C$ can respect the floor $k_{\min}$ \EndIf
\If{$k=2$} \Return $(\{C_1\},\{C_2\})$ \Comment{trivial; matches what the general path below gives}
\EndIf
\State compute $W$
\If{$W\approx 0$} \Comment{no resolvable correlation}
    \State order $\gets C$ in its own given index order; \quad cut $\gets \lfloor k/2\rfloor$
\Else
    \State compute Laplacian $L$ and Fiedler vector $v$
    \If{$\mathrm{sign}(v)$ is the same for every entry} \Comment{degenerate split}
        \State order $\gets C$ sorted by $v$, ascending; \quad cut $\gets \lfloor k/2\rfloor$
    \Else
        \State order $\gets C$ sorted by $v$, ascending; \quad cut $\gets \#\{p: v_p\le 0\}$
    \EndIf
\EndIf
\State cut $\gets \min\!\big(\max(\text{cut},\,k_{\min}),\; k-k_{\min}\big)$ \Comment{clamp; no-op at $k_{\min}=1$}
\State \Return $\big(\text{order}[\text{cut}:],\ \text{order}[:\text{cut}]\big)$
\end{algorithmic}
\end{algorithm}

Orbitals whose occupations fluctuate together end up on the same side: $W_{pq}$ is large exactly
when $n_p$ and $n_q$ are strongly correlated, so cutting along the Fiedler vector separates $C$
into two weakly-correlated groups -- a cheap proxy for the split that will end up costing the
least under the subsequent continuous optimization.

The two fallback branches use different orderings even though both cut at $\lfloor k/2\rfloor$:
the no-correlation branch keeps $C$'s own given order, since there is no correlation structure
to sort by; the degenerate-Fiedler branch sorts by the Fiedler vector itself, since there is
structure, just not one with a clean sign split. Keeping these separate matters for
reproducibility of a given run, not for correctness: neither ordering is more ``right'' than the
other in the degenerate case.

The clamp step is the feature this section documents, and is provably a no-op at the default
$k_{\min}=1$: every branch above already produces two non-empty groups, so the natural cut is
always already in $[1,k-1]$, a subset of the clamp range at $k_{\min}=1$. For $k_{\min}>1$ the
clamp pulls the boundary in from whichever side is too small, moving exactly the
Fiedler-boundary-nearest orbitals across -- not an arbitrary reshuffling. The $k=2$ early return
is not a fourth case: it only skips a trivial $2\times2$ eigendecomposition that would return the
same split. The Laplacian eigendecomposition costs $O(k^3)$ for a cluster of size $k$, negligible
next to the $O(\mathrm{norb}^5)$-or-worse cost evaluations elsewhere: Fiedler splitting is never
the bottleneck.

\section{Beam Search}
\label{sec:beam}

\begin{center}
\begin{tabular}{ll}
\toprule
Symbol & Meaning \\
\midrule
$(P,U)$ & a decomposition: partition plus accumulated rotation \\
$B$ & the beam: a list of up to \texttt{num\_decos} decompositions \\
$b$, $s$ & \texttt{num\_decos} (beam width), \texttt{num\_subdecos} (children/parent/round) \\
$k_p$, $k_c$ & \texttt{min\_parent\_cluster\_size}, \texttt{min\_child\_cluster\_size} \\
$T$ & \texttt{target\_num\_clusters} (optional early-stop target) \\
$\mathrm{traj}$ & the trajectory: best decomposition found for each cluster count \\
\textsc{CanSplit}$(k)$ & $(k>k_p)\wedge(k\ge 2k_c)$: is a cluster of size $k$ worth trying to split \\
\bottomrule
\end{tabular}
\end{center}

\begin{algorithm}[h]
\caption{\textsc{RunDecompositionOptimizer}}
\label{alg:beam}
\begin{algorithmic}[1]
\small
\State $B \gets$ $b$ copies of the trivial 1-cluster decomposition, seeded from reference bases
    (MOs and/or NatOs, cycled, optionally each Fiedler-reordered)
\State $\mathrm{traj} \gets [\arg\min_{\text{cost}} B]$
\While{true}
  \If{$T$ set and every $(P,U)\in B$ has $\abs{P}\ge T$} \textbf{break} \EndIf
  \If{no cluster in any $(P,U)\in B$ satisfies \textsc{CanSplit}} \textbf{break} \EndIf
  \If{round-count cap reached} \textbf{break} \EndIf
  \State candidates $\gets \emptyset$
  \For{each $(P,U) \in B$}
    \State rank $P$'s clusters satisfying \textsc{CanSplit} by current $\Var(N_C)$, descending
    \For{each of the top $s$ such clusters $C$}
      \State $(V_1,V_2) \gets \textsc{FiedlerBipartition}(n_1,n_2,C,k_c)$ \Comment{Algorithm~\ref{alg:fiedler}}
      \State $P' \gets (P\setminus\{C\})\cup\{V_1,V_2\}$ \Comment{only $C$ changes}
      \State re-optimize $U$, restricted to $C$'s internal generator entries, warm-started at $0$
      \State candidates $\gets$ candidates $\cup \{(P',U')\}$
    \EndFor
  \EndFor
  \If{candidates $=\emptyset$} \textbf{break} \Comment{provably unreachable if line 5 did not fire}\EndIf
  \State $B \gets$ best $b$ of candidates, pooled globally across all parents, by cost
  \State $\mathrm{traj} \gets \mathrm{traj} + [\arg\min_{\text{cost}} B]$
\EndWhile
\State \Return $\mathrm{traj}$
\end{algorithmic}
\end{algorithm}

\paragraph{Discrete versus continuous.} \emph{Which} cluster to split is decided by a cheap,
cost-function-independent rule: among clusters satisfying \textsc{CanSplit}, pick those with the
largest current $\Var(N_C)$. Only the continuous re-optimization of $U$ after a split, and the
ranking used to prune the beam, depend on the user-supplied cost function. This is why the cheap
variance cost should drive the search itself: the discrete decisions never depend on which cost
function is passed in, but every candidate's continuous optimization does, so an expensive cost
multiplies its cost by the number of candidates touched.

\paragraph{Block-restricted optimization.} When cluster $C$ is split, only the entries of the
antisymmetric generator $A$ mixing pairs of orbitals \emph{within} $C$ are freed for the
subsequent L-BFGS re-optimization; every other entry of $A$ stays zero. This is \emph{not} a
cheaper gradient evaluation: the implementation always builds the full
$\binom{\mathrm{norb}}{2}$-length, zero-padded parameter vector and calls $\exp(\cdot)$ on a
dense $\mathrm{norb}\times\mathrm{norb}$ matrix regardless of how many entries are free, and
reverse-mode automatic differentiation costs about the same multiple of the forward pass
independent of how many inputs are free. The real gains are a lower-dimensional optimization
landscape (fewer local optima, fewer outer iterations in practice), and, more importantly, an
exact structural guarantee. Because $A$ is zero outside $C$'s block, $\exp(A)$ is
block-diagonal: identity on every orbital outside $C$. Splitting $C$ therefore leaves every
other cluster's orbitals, and their number-operator statistics, exactly unchanged -- the
\emph{block-locality invariant} that makes it meaningful to compare candidate splits of
different clusters against a common, otherwise-frozen partition. Since $A$ is always real and
antisymmetric, $U=\exp(A)\in SO(\mathrm{norb})$ always, and $SO(\mathrm{norb})$ is closed under
composition, so every accumulated rotation stays a proper rotation throughout the search.

\paragraph{The search loop.} Each child differs from its parent in exactly one cluster: the one
chosen for splitting. Every round, every beam member proposes up to $s$ children; all children
from all beam members, from every parent, are pooled into one global candidate list, and the
best $b$ by cost become the new beam -- not the best $s$ per parent kept separately. This global
pooling means a single strong parent can, in principle, supply all $b$ survivors, at the cost of
losing lineage diversity, but it also guarantees the beam's best cost is non-increasing round
over round. Because every surviving beam member gains exactly one cluster per round, starting
from a uniform 1-cluster beam, every beam member has the same cluster count after any given
number of rounds; by induction, the trajectory has no gaps and no duplicates:
$\mathrm{traj}[i]$ has exactly $i+1$ clusters. The restricted re-optimization is warm-started at
the all-zero parameter vector -- the raw Fiedler cut, with no further rotation -- which is a
valid point of the restricted search space, so the optimized cost is guaranteed no worse than the
cost at that starting point.

\paragraph{Stopping and the trajectory.} The loop stops when (1) a requested target cluster
count is reached; (2) no cluster anywhere in the beam satisfies \textsc{CanSplit}; (3) a
round-count safety cap is hit; or (4), defensively, no children were produced in a round.
Condition (4) is provably redundant with (2): a beam member contributes at least one child
whenever it has a splittable cluster, so if (2) had not already fired, (4) cannot fire either --
it is kept only for defensive symmetry against future changes to the eligibility rule. The
result is a \emph{trajectory}: the best decomposition found for each cluster count reached, not
a single winner. This is necessary because raw cost, on its own, trivially favors finer
partitions. The clearest instance is an exact operator identity, not an empirical observation:
for the un-split, single all-orbitals cluster, $\Var(N_\mathrm{total})=0$ for \emph{every}
rotation $U$, because total particle number is invariant under any single-particle basis change
and the state has fixed particle number -- no optimization is needed to see this. Picking a
useful point on the trajectory therefore requires an external metric;
Section~\ref{sec:sectors}'s convergence analysis is that metric in practice, and
\texttt{best\_decomposition(traj, num\_clusters=$k$)} retrieves the trajectory entry with
exactly $k$ clusters once a preferred $k$ is known. Every candidate partition, at every round,
is validated -- covers every orbital exactly once, no empty clusters -- before it is accepted
into the beam.

\section{Polishing}
\label{sec:polish}

After the beam search terminates, each trajectory entry can optionally undergo one further, full
joint reoptimization: all $\binom{\mathrm{norb}}{2}$ generator entries are freed at once (not
block-restricted), starting from the beam search's own $U$, for that entry's now-fixed
partition. This polishes away suboptimality the greedy, block-restricted process could not reach
on its own. Polishing runs once per trajectory entry, entirely independently of the others: it
does not feed back into how the trajectory was built. As with each split, it is warm-started at
the beam search's own rotation, so the polished cost is guaranteed no worse than the unpolished
one -- and it is the natural place to switch to a more expensive, more physically faithful cost
such as the commutator cost, since it now runs only once per entry rather than once per
beam-search candidate.

\section{Downstream Sector Analysis}
\label{sec:sectors}

Selected trajectory entries -- by default, every entry with at least two clusters, since a
single all-orbitals cluster carries no symmetry information -- are passed through the same
K-sectors/K-states Hilbert-space-truncation convergence analysis used by
\texttt{cluster\_numbers\_metrics.py}'s fixed-partition workflow, reused essentially unchanged:
for each entry and each requested maximum electron-transfer count, it ranks symmetry sectors
(respectively, decoupled eigenstates) by proximity to the reference state, and reports how the
projected energy and retained dimension converge as more are kept, together with whether
chemical accuracy is reached. The only structural difference is bookkeeping: there, one
\texttt{cluster\_matrix} is shared across several bases, so sectors are computed once; here,
every entry carries its own partition, so sectors are rebuilt per entry. This is the practical
answer to Section~\ref{sec:beam}'s open question of which trajectory entry to prefer: the one
reaching chemical accuracy with fewest retained sectors or states, not the one with lowest raw
cost.

\section{Usage}

\begin{Verbatim}[fontsize=\footnotesize]
python cluster_number_decomposition_optimization.py h2o sto-3g 2.0 variance \
    --bond-angle 104.5 --initial-basis both --target-num-clusters 3
\end{Verbatim}

\begin{center}
\begin{tabularx}{\textwidth}{p{4.3cm}X}
\toprule
Flag & Meaning \\
\midrule
\texttt{--num-decos}, \texttt{--num-subdecos} & beam width $b$, children per parent per round $s$ (default 4, 4) \\
\texttt{--min-parent-cluster-size} & $k_p$: clusters at or below this size are never chosen for splitting (default 1) \\
\texttt{--min-child-cluster-size} & $k_c$: neither child of a split may end up smaller than this (default 1) \\
\texttt{--target-num-clusters} & stop once this many clusters is reached (default: run to singletons) \\
\texttt{--initial-basis \{MOs,NatOs,both\}} & reference basis (or bases) seeding the beam (default MOs) \\
\texttt{--fiedler-reorder} & reorder each initial basis by its own global Fiedler vector \\
\texttt{--no-polish}, \texttt{--polish-maxiter} & skip, or tune the iteration budget of, Section~\ref{sec:polish}'s polish step \\
\texttt{--analyze-num-clusters}, \texttt{--skip-sector-analysis}, \texttt{--skip-K-sectors}, \texttt{--skip-K-states} & which trajectory entries get Section~\ref{sec:sectors}'s analysis, and which half of it \\
\bottomrule
\end{tabularx}
\end{center}

DMRG, HPC, logging and output-path flags (\texttt{--bond-dim}, \texttt{--n-sweeps},
\texttt{--n-threads}, \texttt{--output-dir}, \dots) are shared verbatim with
\texttt{cluster\_numbers\_metrics.py}; see \texttt{--help} for the full list.

As a library, the search can be driven directly on RDMs computed elsewhere:

\begin{Verbatim}[fontsize=\footnotesize]
from cluster_number_decomposition_optimization import (
    RDMData, DecompositionOptimizerConfig, run_decomposition_optimizer,
    make_variance_cost_constructor, polish_decomposition, best_decomposition,
)
rdm_data = RDMData(D=D, Gamma=Gamma)
config = DecompositionOptimizerConfig(num_decos=4, num_subdecos=6, min_parent_cluster_size=2)
traj = run_decomposition_optimizer(make_variance_cost_constructor(), rdm_data, config)
deco = best_decomposition(traj, num_clusters=3)
deco = polish_decomposition(deco, rdm_data, make_variance_cost_constructor())
\end{Verbatim}

\end{document}
